\section*{Bab \ref{ch:g4:clouds-rain-snow} --- Awan, Hujan, dan Salju}

\begin{solution}{exo:g4:clouds-rain-snow:1}
\omterm{def:g4:clouds-rain-snow:cloud}{Awan} tersusun dari sejumlah besar titik air yang sangat kecil, atau
kepingan es yang sangat kecil, yang melayang di udara. \omterm{def:g4:clouds-rain-snow:cloud}{Kabut} adalah \omterm{def:g4:clouds-rain-snow:cloud}{awan}
yang menyentuh tanah.
\end{solution}

\begin{solution}{exo:g4:clouds-rain-snow:2}
(a) \omterm{def:g4:clouds-rain-snow:cloud-families}{Stratus}. (b) \omterm{def:g4:clouds-rain-snow:cloud-families}{Sirus}. (c) \omterm{def:g4:clouds-rain-snow:cloud-families}{Kumulus}. (d) \omterm{def:g4:clouds-rain-snow:cloud-families}{Kumulonimbus}.
\end{solution}

\begin{solution}{exo:g4:clouds-rain-snow:3}
(a) \omterm{def:g4:clouds-rain-snow:precipitation}{Hujan es}. (b) \omterm{def:g4:clouds-rain-snow:precipitation}{Hujan}. (c) \omterm{def:g4:clouds-rain-snow:precipitation}{Salju}.
\end{solution}

\begin{solution}{exo:g4:clouds-rain-snow:4}
(a) \omterm{def:g4:clouds-rain-snow:rain-gauge}{Penakar hujan}. (b) Termometer. (c) \omterm{def:g4:clouds-rain-snow:wind-vane}{Penunjuk arah angin}.
\end{solution}

\begin{solution}{exo:g4:clouds-rain-snow:5}
\omterm{def:g4:clouds-rain-snow:precipitation}{Hujan} yang sudah turun 17 milimeter. $20 - 17 = 3$: 3 milimeter lagi akan
membawa air sampai ke tanda 20.
\end{solution}

\begin{solution}{exo:g4:clouds-rain-snow:6}
Hari yang paling basah adalah Rabu (12 milimeter). Tidak turun \omterm{def:g4:clouds-rain-snow:precipitation}{hujan} pada
hari Senin dan Jumat. Hari Kamis mendapat 7 milimeter.
\end{solution}

\begin{solution}{exo:g4:clouds-rain-snow:7}
Kepingan \omterm{def:g4:clouds-rain-snow:precipitation}{salju} terbentuk tinggi di dalam \omterm{def:g4:clouds-rain-snow:cloud}{awan} yang dingin. Di desa
pegunungan di atas, udara membeku sepanjang jalan turun, jadi kepingan itu
tiba sebagai \omterm{def:g4:clouds-rain-snow:precipitation}{salju}. Lebih ke bawah, di \omterm{def:g1:landscapes-around-us:landform}{lembah}, udaranya lebih hangat
daripada \SI{0}{\celsius}: dalam perjalanan turun melewati udara itu,
kepingan \omterm{def:g4:clouds-rain-snow:precipitation}{salju} mencair, dan sampai ke tanah sebagai \omterm{def:g4:clouds-rain-snow:precipitation}{hujan}.
\end{solution}

\begin{solution}{exo:g4:clouds-rain-snow:8}
Diagram batangnya mempunyai tujuh batang dengan tinggi 3, 0, 1, 9, 6, 0,
dan 1 (tidak ada batang pada hari Selasa dan Sabtu). Dalam minggu itu,
turun $3 + 0 + 1 + 9 + 6 + 0 + 1 = 20$ milimeter. Hari Kamis membawa 9
dari 20 milimeter itu: $\frac{9}{20}$ dari curah \omterm{def:g4:clouds-rain-snow:precipitation}{hujan} minggu itu.
\end{solution}

\begin{solution}{exo:g4:clouds-rain-snow:9}
Dinding dan pohon melindungi tanah di dekatnya dari \omterm{def:g4:clouds-rain-snow:precipitation}{hujan}, terutama ketika
angin meniup \omterm{def:g4:clouds-rain-snow:precipitation}{hujan} ke samping, dan tetes air bisa menetes dari daun ke
dalam penakar: di dekatnya, penakar akan menampung terlalu sedikit, atau
terlalu banyak. Termometer harus mengukur suhu udara; di bawah terik
matahari, termometer akan dihangatkan oleh sinar matahari itu sendiri dan
menunjukkan angka yang terlalu tinggi. Kotak putih berkisi-kisi menahan
sinar matahari dan membiarkan udara lewat.
\end{solution}

\begin{solution}{exo:g4:clouds-rain-snow:10}
\omterm{def:g4:clouds-rain-snow:wind-vane}{Penunjuk arah angin} menunjuk ke arah asal angin: angin datang dari barat.
Angin itu pergi ke arah timur.
\end{solution}

\begin{solution}{exo:g4:clouds-rain-snow:11}
\omterm{def:g4:clouds-rain-snow:precipitation}{Hujan} satu milimeter di atas satu meter persegi sama dengan satu liter. Di
atas satu meter persegi, 8 milimeter memberi 8 liter. Di atas 500 meter
persegi: $500 \times 8 = 4000$ liter. Dengan gembor berisi 10 liter:
$4000 \div 10 = 400$ gembor.
\end{solution}
